% Layered hyperedge rendering: graph (a) + computational (c)
% Same 4-vertex / 3-hyperedge data; two complementary views.
% Used by closure_quotient_with_hodge_corpus.tex §1.7.
\begin{tikzpicture}[font=\small]

% ---------- Panel (a): graph view (centroid-star) ----------
\begin{scope}[shift={(0,0)}]
  \node[font=\bfseries] at (0,3.0) {(a) Graph view};
  \node[font=\scriptsize,align=center,text width=4.6cm] at (0,2.55)
       {ternary hyperedge = three vertices linked by a centroid};

  \coordinate (av1) at (0, 1.5);
  \coordinate (av2) at (-1.3, 0);
  \coordinate (av3) at (0, -1.5);
  \coordinate (av4) at (1.3, 0);

  \coordinate (ae1) at ({(0 + (-1.3) + 0)/3}, {(1.5 + 0 + (-1.5))/3});
  \coordinate (ae2) at ({((-1.3) + 0 + 1.3)/3}, {(0 + (-1.5) + 0)/3});
  \coordinate (ae3) at ({(0 + 0 + 1.3)/3}, {(1.5 + (-1.5) + 0)/3});

  \draw[blue!60,thin] (av1)--(ae1)--(av2) (ae1)--(av3);
  \draw[red!60,thin]  (av2)--(ae2)--(av3) (ae2)--(av4);
  \draw[green!50!black,thin] (av1)--(ae3)--(av3) (ae3)--(av4);

  \fill[blue!60]  (ae1) circle (1.3pt);
  \fill[red!60]   (ae2) circle (1.3pt);
  \fill[green!50!black] (ae3) circle (1.3pt);

  \foreach \v/\lbl in {av1/$v_1$, av2/$v_2$, av3/$v_3$, av4/$v_4$} {
    \fill[black] (\v) circle (1.6pt);
    \node[font=\scriptsize, above right=-1pt and 1pt of \v] {\lbl};
  }
\end{scope}

% ---------- Panel (b): computational view (ordered-slot) ----------
\begin{scope}[shift={(7.5,0)}]
  \node[font=\bfseries] at (0,3.0) {(b) Computational view};
  \node[font=\scriptsize,align=center,text width=4.6cm] at (0,2.55)
       {slot 1, 2 = operator inputs; slot 3 = output (arrow)};

  \coordinate (cv1) at (0, 1.5);
  \coordinate (cv2) at (-1.3, 0);
  \coordinate (cv3) at (0, -1.5);
  \coordinate (cv4) at (1.3, 0);

  \coordinate (ce1) at ({(0 + (-1.3) + 0)/3}, {(1.5 + 0 + (-1.5))/3});
  \coordinate (ce2) at ({((-1.3) + 0 + 1.3)/3}, {(0 + (-1.5) + 0)/3});
  \coordinate (ce3) at ({(0 + 0 + 1.3)/3}, {(1.5 + (-1.5) + 0)/3});

  % e1 = (v1, v2, v3): inputs solid, output dashed
  \draw[blue!60,thick] (cv1)--(ce1);
  \draw[blue!60,thick] (cv2)--(ce1);
  \draw[blue!60,dashed,-{Latex[length=4pt]}] (ce1)--(cv3);

  \draw[red!60,thick] (cv2)--(ce2);
  \draw[red!60,thick] (cv3)--(ce2);
  \draw[red!60,dashed,-{Latex[length=4pt]}] (ce2)--(cv4);

  \draw[green!50!black,thick] (cv1)--(ce3);
  \draw[green!50!black,thick] (cv3)--(ce3);
  \draw[green!50!black,dashed,-{Latex[length=4pt]}] (ce3)--(cv4);

  \fill[blue!60]  (ce1) circle (1.3pt);
  \fill[red!60]   (ce2) circle (1.3pt);
  \fill[green!50!black] (ce3) circle (1.3pt);

  \foreach \v/\lbl in {cv1/$v_1$, cv2/$v_2$, cv3/$v_3$, cv4/$v_4$} {
    \fill[black] (\v) circle (1.6pt);
    \node[font=\scriptsize, above right=-1pt and 1pt of \v] {\lbl};
  }
\end{scope}

\end{tikzpicture}
